\documentclass[12pt, a4paper]{report}

% ─── Packages ────────────────────────────────────────────────────────────────
\usepackage[utf8]{inputenc}
\usepackage[T1]{fontenc}
\usepackage{lmodern}
\usepackage[english]{babel}
\usepackage{geometry}
\usepackage{graphicx}
\usepackage{amsmath, amssymb}
\usepackage{booktabs}
\usepackage{array}
\usepackage{hyperref}
\usepackage{setspace}
\usepackage{titlesec}
\usepackage{fancyhdr}
\usepackage{tocloft}
\usepackage{xcolor}
\usepackage{listings}
\usepackage{caption}
\usepackage{subcaption}
\usepackage{enumitem}
\usepackage{parskip}
\usepackage{float}
\usepackage{cite}
\usepackage{algorithm}
\usepackage{algpseudocode}

% ─── Page geometry ───────────────────────────────────────────────────────────
\geometry{
  top=2.5cm, bottom=2.5cm,
  left=3cm,  right=2.5cm
}

% ─── Spacing ─────────────────────────────────────────────────────────────────
\onehalfspacing

% ─── Colors ──────────────────────────────────────────────────────────────────
\definecolor{vnuBlue}{RGB}{0, 70, 127}
\definecolor{accentRed}{RGB}{180, 30, 30}
\definecolor{codeGray}{RGB}{245, 245, 245}
\definecolor{commentGreen}{RGB}{60, 128, 49}

% ─── Hyperref setup ──────────────────────────────────────────────────────────
\hypersetup{
  colorlinks = true,
  linkcolor  = vnuBlue,
  citecolor  = accentRed,
  urlcolor   = vnuBlue,
  pdftitle   = {Video Summarization by Supervised Learning},
  pdfauthor  = {Tran Hai Duc}
}

% ─── Header / Footer ─────────────────────────────────────────────────────────
\pagestyle{fancy}
\fancyhf{}
\fancyhead[L]{\footnotesize\textit{Video Summarization by Supervised Learning}}
\fancyhead[R]{\footnotesize\textit{Duc (23127173) \& Chien (23127331)}}
\fancyfoot[C]{\thepage}
\renewcommand{\headrulewidth}{0.4pt}

% ─── Section title formatting ─────────────────────────────────────────────────
\titleformat{\chapter}[display]
  {\normalfont\Large\bfseries\color{vnuBlue}}
  {\chaptertitlename\ \thechapter}{10pt}{\huge}
\titlespacing*{\chapter}{0pt}{30pt}{20pt}

\titleformat{\section}
  {\normalfont\large\bfseries\color{vnuBlue}}
  {\thesection}{1em}{}

\titleformat{\subsection}
  {\normalfont\normalsize\bfseries}
  {\thesubsection}{1em}{}

% ─── Code listing style ───────────────────────────────────────────────────────
\lstset{
  backgroundcolor=\color{codeGray},
  basicstyle=\small\ttfamily,
  breaklines=true,
  commentstyle=\color{commentGreen},
  keywordstyle=\bfseries\color{vnuBlue},
  numbers=left,
  numberstyle=\tiny\color{gray},
  frame=single,
  captionpos=b
}

% ─────────────────────────────────────────────────────────────────────────────
%  DOCUMENT
% ─────────────────────────────────────────────────────────────────────────────
\begin{document}

% ════════════════════════════════════════════════════════════════════════════
%  COVER PAGE
% ════════════════════════════════════════════════════════════════════════════
\begin{titlepage}
\thispagestyle{empty}
\begin{center}

  {\large\textbf{VIETNAM NATIONAL UNIVERSITY -- HO CHI MINH CITY}}\\[4pt]
  {\large\textbf{UNIVERSITY OF SCIENCE}}\\[4pt]
  {\large\textbf{FACULTY OF INFORMATION TECHNOLOGY}}

  \vspace{1.5cm}
  \hrule height 1.2pt
  \vspace{0.4cm}
  \hrule height 0.4pt
  \vspace{1.5cm}

  % \includegraphics[width=3cm]{hcmus_logo.png}\\[1cm]

  {\color{vnuBlue}\Huge\bfseries
    Video Summarization\\[8pt]
    by Supervised Learning
  }

  \vspace{0.8cm}
  {\Large\bfseries Research Report}

  \vspace{0.5cm}
  \hrule height 0.4pt
  \vspace{0.4cm}
  \hrule height 1.2pt
  \vspace{1.8cm}

  \begin{minipage}{0.48\textwidth}
    \raggedright
    \textbf{Students:}\\[4pt]
    \quad Tran Hai Duc\\
    \quad Student ID: \texttt{23127173}\\[6pt]
    \quad Nguyen Cong Chien\\
    \quad Student ID: \texttt{23127331}
  \end{minipage}
  \hfill
  \begin{minipage}{0.48\textwidth}
    \raggedleft
    \textbf{Supervisor:}\\[4pt]
    Vo Hoai Viet, Ph.D.
  \end{minipage}

  \vfill

  {\large Ho Chi Minh City, \the\year}

\end{center}
\end{titlepage}

% ════════════════════════════════════════════════════════════════════════════
%  ACKNOWLEDGEMENTS
% ════════════════════════════════════════════════════════════════════════════
\chapter*{Acknowledgements}
\addcontentsline{toc}{chapter}{Acknowledgements}
\thispagestyle{fancy}

First and foremost, I would like to express my sincere gratitude to my supervisor,
\textbf{Dr.\ Vo Hoai Viet}, for his invaluable guidance, encouragement, and continuous
support throughout this research project. His insightful feedback, deep expertise in
computer vision and machine learning, and willingness to create a conducive research
environment have been instrumental in shaping both the direction and the quality of
this work. I am truly grateful for every discussion, suggestion, and resource he
generously provided.

I also wish to extend my heartfelt thanks to the research community whose open
contributions have made this study possible. The authors of the \textbf{SumMe} and
\textbf{TVSum} benchmark datasets provided critical evaluation resources that enabled
rigorous experimentation. The developers of \textbf{OpenAI Whisper}, \textbf{SentenceBERT},
and the underlying pre-trained deep learning models offered foundational building blocks
without which the proposed pipeline could not have been realised. I likewise acknowledge
the pioneering work of \textbf{Zhang et al.}\ on video summarization with long
short-term memory networks, which directly inspired the architectural choices in this
project.

Special thanks go to my teammate \textbf{Nguyen Cong Chien} (Student ID: 23127331)
for his dedication, collaboration, and shared effort throughout every stage of this
project -- from system design and implementation to experimentation and writing.
I also thank my classmates and friends for their moral support and stimulating
discussions during challenging moments of the research process.

Finally, I am deeply indebted to my family for their unwavering love, patience, and
encouragement throughout my studies.

This work was carried out at the University of Science, Vietnam National University --
Ho Chi Minh City. Any remaining errors or shortcomings are solely my own responsibility.

\vspace{1.5cm}
\begin{flushright}
  \textit{Ho Chi Minh City, \the\year}\\[6pt]
  \textit{Tran Hai Duc}
\end{flushright}

\newpage

% ════════════════════════════════════════════════════════════════════════════
%  TABLE OF CONTENTS
% ════════════════════════════════════════════════════════════════════════════
\tableofcontents
\newpage

% ════════════════════════════════════════════════════════════════════════════
%  CHAPTER 1 -- INTRODUCTION
% ════════════════════════════════════════════════════════════════════════════
\chapter{Introduction}

\section{Overview of Video Summarization}

\textbf{Video summarization} is the task of automatically generating a compact
representation of a video that retains its most important content and storyline
\cite{haq2021,money2015}. Depending on the output format, summaries fall into two
main categories:

\begin{itemize}[leftmargin=1.5cm]
  \item \textbf{Static summary:} a set of representative keyframes extracted from
        the original video.
  \item \textbf{Dynamic summary:} a shortened video clip assembled from selected
        important segments, preserving both visual content and audio.
\end{itemize}

The overarching goal is to \emph{reduce viewing time} without sacrificing the essential
narrative, allowing users to grasp the key content efficiently.

\section{Motivating Examples}

\subsection{Everyday Use Case}
Consider a user who wishes to watch a one-hour lecture on YouTube but has only ten
minutes available. A video summarization system can:
\begin{enumerate}
  \item Automatically identify segments containing core explanations, motivating
        examples, and concluding remarks.
  \item Assemble a condensed version of approximately 6--10 minutes, enabling the user
        to absorb 70--80\,\% of the essential content without watching the full hour.
\end{enumerate}

\subsection{Data Reduction Perspective}
In \textbf{video classification} tasks, models must process thousands of frames per
video, incurring large computational costs. By first summarizing a video to retain
only $\approx 15\,\%$ of its frames, the downstream classifier effectively operates on
a lower-dimensional representation. This yields two concrete benefits:
\begin{itemize}
  \item Reduced computation and storage requirements.
  \item A focus on information-rich frames, which can improve classification accuracy
        on the same hardware budget.
\end{itemize}

\section{Significance of the Work}

\subsection{Practical Significance}
The practical impact of video summarization includes \cite{kadam2022}:
\begin{itemize}
  \item \textbf{Time savings for users:} viewers can quickly understand the content
        of long videos without watching them in full.
  \item \textbf{Support for recommendation and retrieval:} meaningful keyframes improve
        thumbnail generation and enable preview functionality.
  \item \textbf{Applications in education and news:} automatic highlight extraction
        for lectures, MOOCs, and broadcast news programmes.
\end{itemize}

\subsection{Scientific Significance}
From a data-science perspective:
\begin{itemize}
  \item \textbf{Structured dimensionality reduction:} video summarization selects
        frames or segments with high \emph{information content}, rather than applying
        uniform down-sampling.
  \item \textbf{Pre-processing for downstream models:} summarizing large video corpora
        before training classification, retrieval, or captioning models reduces
        training costs and focuses learning on informative data.
\end{itemize}

\section{Research Context}

Current approaches to video summarization span a wide spectrum \cite{apostolidis2021}:
\begin{itemize}
  \item \textbf{Traditional methods:} heuristic rules, clustering, shot detection.
  \item \textbf{Deep learning:} sequence labeling with LSTM/BiLSTM \cite{zhang2016},
        attention and Transformer models, reinforcement learning, self-supervised
        learning.
  \item \textbf{Multimodal methods:} combining visual, audio, transcript, and metadata
        features \cite{psallidas2021,hu2023}.
\end{itemize}

Recent surveys \cite{apostolidis2021} identify several open opportunities:
\begin{itemize}
  \item Better exploitation of textual and audio signals (speech, subtitles, OCR).
  \item More effective modeling of long sequences using hierarchical Transformers.
  \item Unsupervised and semi-supervised learning to reduce annotation costs.
\end{itemize}

\section{Research Focus}

This project targets \textbf{multimodal video summarization (visual + audio)} in a
supervised setting on the SumMe \cite{gygli2014} and TVSum \cite{song2015} benchmarks.
The proposed architecture combines a \textbf{BiLSTM with Temporal Attention}
\cite{zhang2016} with an audio branch built on Whisper \cite{radford2023} and
SentenceBERT \cite{reimers2019}. The primary deliverable is a \textbf{dynamic summary}
-- a condensed video clip with intact audio.

% ════════════════════════════════════════════════════════════════════════════
%  CHAPTER 2 -- PROBLEM FORMULATION
% ════════════════════════════════════════════════════════════════════════════
\chapter{Problem Formulation}

\section{Input}

\begin{itemize}
  \item A raw video $V$ in any standard container format (\texttt{.mp4}, \texttt{.avi},
        etc.) of arbitrary duration.
  \item During training: frame-level importance scores $y_t \in [0,1]$ for each frame
        $t$, derived from human annotations provided in the SumMe and TVSum datasets.
\end{itemize}

\section{Output}

\begin{itemize}
  \item \textbf{Static summary:} a set of representative keyframes that collectively
        depict the video's main content.
  \item \textbf{Dynamic summary:} a condensed video $V_{\text{sum}}$ of approximately
        $15\,\%$ of the original duration, assembled from the most important segments
        and retaining the corresponding audio track.
\end{itemize}

\section{Implementation Focus}

This project concentrates on \textbf{dynamic summaries} for three reasons:
\begin{enumerate}
  \item They more closely match real-world user needs -- watching a short video with
        sound rather than browsing static images.
  \item They allow qualitative evaluation through direct viewing experience.
  \item They fully exploit the \texttt{ffmpeg}-based audio processing pipeline
        developed in this project.
\end{enumerate}

Static keyframe export is also supported as a by-product of the keyshot selection step.

% ════════════════════════════════════════════════════════════════════════════
%  CHAPTER 3 -- METHODOLOGY
% ════════════════════════════════════════════════════════════════════════════
\chapter{Methodology}

\section{Supervised Learning Framework}

The model learns to approximate human importance judgements:
\begin{itemize}
  \item \textbf{Input:} a sequence of multimodal feature vectors
        $\{x_t\}_{t=1}^{T}$, $x_t \in \mathbb{R}^{1408}$.
  \item \textbf{Output:} predicted importance scores $\hat{y}_t \in [0,1]$ for each
        frame $t$.
  \item \textbf{Ground truth:} $y_t \in [0,1]$, the mean score assigned by 15--20
        human annotators.
  \item \textbf{Loss:} L1 (or MSE) loss between $\hat{y}_t$ and $y_t$.
\end{itemize}

\section{Datasets}

\subsection{SumMe \cite{gygli2014}}
SumMe is a collection of short, thematically diverse user-generated videos.
Annotators manually mark important segments; these are converted to frame-level
importance scores used as ground truth.

\subsection{TVSum \cite{song2015}}
TVSum is oriented toward ``TV program'' content (news, documentaries, vlogs). Each
video is annotated by approximately 20 raters; the mean score constitutes the ground
truth.

\section{Evaluation Metrics}

\begin{itemize}
  \item \textbf{F-score} (F-measure): measures temporal overlap between the generated
        summary and the annotator-derived reference summary.
  \item \textbf{Precision / Recall:}
        \begin{itemize}
          \item \emph{Precision} -- proportion of summary content that is genuinely important.
          \item \emph{Recall} -- proportion of important video content retained in the summary.
        \end{itemize}
  \item \textbf{Temporal Overlap:} degree of time-segment overlap between predicted
        summary segments and ground-truth segments.
\end{itemize}

% ════════════════════════════════════════════════════════════════════════════
%  CHAPTER 4 -- PROPOSED ARCHITECTURE
% ════════════════════════════════════════════════════════════════════════════
\chapter{Proposed Architecture: Multimodal BiLSTM with Temporal Attention}

\section{BiLSTM Overview}

A \textbf{Bidirectional Long Short-Term Memory (BiLSTM)} network \cite{zhang2016}
processes the input sequence in both the forward (past) and backward (future) temporal
directions. For each frame $t$, the hidden state encodes:
\begin{itemize}
  \item Context from all preceding frames (forward pass).
  \item Context from all subsequent frames (backward pass).
\end{itemize}

Formally, given input sequence $(x_1, \dots, x_T)$ with $x_t \in \mathbb{R}^{1408}$:
\begin{equation}
  (h_1, \dots, h_T) = \mathrm{BiLSTM}(x_1, \dots, x_T), \quad h_t \in \mathbb{R}^{2H}
\end{equation}
where $H$ is the hidden size per direction (set to $256$, giving
$h_t \in \mathbb{R}^{512}$).

\section{Temporal Attention}

Temporal attention is applied on top of the BiLSTM outputs to focus the model on
the most informative time steps:
\begin{align}
  e_t      &= \tanh\!\bigl(W_h\, h_t + b_h\bigr) \\
  \alpha_t &= \operatorname{softmax}\!\bigl(w^{\top} e_t\bigr)
\end{align}
The weights $\alpha_t$ modulate each frame's contribution when predicting importance
scores, effectively down-weighting noisy or redundant segments.

\section{Frame Importance Prediction}

Each hidden state is passed through a linear output layer with sigmoid activation:
\begin{equation}
  \hat{y}_t = \sigma\!\bigl(W_o\, h_t + b_o\bigr), \quad \hat{y}_t \in [0,1]
\end{equation}
producing the score sequence $\{\hat{y}_t\}_{t=1}^{T}$.

\section{Motivation for BiLSTM}

\begin{itemize}
  \item \textbf{Long-range context:} frame importance depends on temporal context,
        not on local content alone.
  \item \textbf{Bidirectionality:} combining past and future context yields richer
        representations than a unidirectional LSTM.
  \item \textbf{Synergy with attention:} attention weights highlight temporally
        important regions while suppressing irrelevant ones.
\end{itemize}

% ════════════════════════════════════════════════════════════════════════════
%  CHAPTER 5 -- AUDIO PROCESSING PIPELINE
% ════════════════════════════════════════════════════════════════════════════
\chapter{Audio Processing Pipeline}

\section{Overview}

To go beyond visual features, this project incorporates \textbf{speech-derived
semantic embeddings} as an additional modality. The pipeline transforms the raw
audio track into a per-frame feature vector aligned with the visual timeline.

\section{Step-by-Step Pipeline}

\subsection{Step 1 -- Audio Extraction}
The audio track is extracted from the source video using \texttt{ffmpeg}, producing
a \texttt{.wav} file.

\subsection{Step 2 -- Automatic Speech Recognition with Whisper}
OpenAI Whisper \cite{radford2023} transcribes the audio into text segments, each
accompanied by start and end timestamps:
\begin{equation*}
  \text{audio} \;\xrightarrow{\;\text{Whisper}\;}\;
  \bigl\{(\text{sentence}_k,\; t^{\text{start}}_k,\; t^{\text{end}}_k)\bigr\}_{k=1}^{K}
\end{equation*}

\subsection{Step 3 -- Sentence Encoding with SentenceBERT}
Each transcript sentence is encoded into a dense semantic vector using
SentenceBERT \cite{reimers2019}:
\begin{equation}
  a^{\text{text}}_k = \mathrm{SentenceBERT}(\text{sentence}_k) \in \mathbb{R}^{384}
\end{equation}

\subsection{Step 4 -- Frame-Level Alignment}
Given the timestamp of each sampled frame $t$, the audio embedding is assigned as:
\begin{equation}
  a_t =
  \begin{cases}
    a^{\text{text}}_k & \text{if frame } t \text{ falls within }
                          [t^{\text{start}}_k,\; t^{\text{end}}_k] \\[4pt]
    \mathbf{0}        & \text{otherwise (silence or no speech)}
  \end{cases}
\end{equation}
This produces an audio feature sequence of shape $[T, 384]$, temporally aligned
with the visual feature sequence $[T, 1024]$.

\subsection{Step 5 -- Multimodal Feature Fusion}
Visual and audio features are concatenated at the frame level:
\begin{equation}
  x_t = \bigl[v_t;\; a_t\bigr] \in \mathbb{R}^{1408}
\end{equation}
The resulting sequence $\{x_t\}_{t=1}^{T}$ is the sole input to the BiLSTM.

\section{Benefits of the Audio Branch}
\begin{itemize}
  \item Enables the model to detect semantically rich spoken content (definitions,
        conclusions, slide titles).
  \item Complements visual features in low-motion segments where important verbal
        explanation occurs over a static slide.
\end{itemize}

% ════════════════════════════════════════════════════════════════════════════
%  CHAPTER 6 -- TRAINING AND INFERENCE ALGORITHMS
% ════════════════════════════════════════════════════════════════════════════
\chapter{Training and Inference Algorithms}

\section{Training Algorithm}

\subsection{Phase 1 -- Data Preprocessing and Feature Extraction}

\paragraph{Visual Features.}
Frames are sampled at 2\,fps using a \texttt{FrameSampler}. A pre-trained CNN
(GoogLeNet/ResNet) extracts a feature vector for each frame:
\begin{equation}
  v_t = \mathrm{CNN}(f_t) \in \mathbb{R}^{1024}
\end{equation}
Features are stored as \texttt{\{video\_id\}.npy} with shape $[T, 1024]$.

Audio features of shape $[T, 384]$ are stored as
\texttt{\{video\_id\}\_audio.npy}, following the pipeline of Chapter~5.

\subsection{Phase 2 -- Input Construction and Batching}

For each frame $t$: $x_t = [v_t;\, a_t] \in \mathbb{R}^{1408}$.
Videos longer than \texttt{max\_seq\_len} (e.g., 960 frames) are chunked.
A \texttt{DataLoader} shuffles videos/segments, pads sequences within each
mini-batch, and applies loss masking on padding positions.

\subsection{Phase 3 -- Forward Pass}

The sequence $\{x_t\}$ is processed by:
\begin{enumerate}
  \item \textbf{BiLSTM encoding} (2 layers, hidden size 256 per direction).
  \item \textbf{Temporal Attention} to compute per-frame weights $\alpha_t$.
  \item \textbf{Output projection} to obtain $\hat{y}_t \in [0,1]$.
\end{enumerate}

\subsection{Phase 4 -- Loss Computation and Optimisation}

\paragraph{Supervised Loss.}
\begin{equation}
  \mathcal{L} = \frac{1}{T} \sum_{t=1}^{T} \bigl|\hat{y}_t - y_t\bigr|
  \quad \text{(L1 loss)}
\end{equation}

\paragraph{Optimiser.}
Adam with learning rate $\eta = 0.001$ (configurable via \texttt{config.yaml}).
Gradients are back-propagated through the BiLSTM, attention, and output layers.

\paragraph{Validation and Early Stopping.}
After each epoch, the validation loss is computed. The best checkpoint is saved to
\texttt{checkpoints/best.pt}. Training halts if no improvement is observed for
\texttt{early\_stopping\_patience} consecutive epochs (triggered at epoch~15 in
the reported experiment).

\section{Inference Algorithm}

\begin{algorithm}[H]
\caption{Dynamic Video Summary Generation}
\begin{algorithmic}[1]
  \Require Raw video $V$, model weights \texttt{best.pt}, summary ratio $r$
  \Ensure Dynamic summary $V_{\text{sum}}$
  \State Extract $\{v_t\}$ and $\{a_t\}$; form $x_t = [v_t; a_t]$
  \State Load model; forward pass $\Rightarrow \{\hat{y}_t\}_{t=1}^{T}$
  \State Apply Gaussian smoothing to $\{\hat{y}_t\}$; normalise
  \State Segment video into contiguous shots
  \State Score each shot (mean of $\hat{y}_t$ within the shot)
  \State Select highest-scoring shots: total duration $\approx r \cdot |V|$,
         with temporal diversity and minimum inter-keyframe gap enforced
  \State Expand each selected segment boundary by $\pm 1.5$\,s
  \State Cut and concatenate segments with audio via \texttt{ffmpeg}
         $\Rightarrow V_{\text{sum}}$
\end{algorithmic}
\end{algorithm}

% ════════════════════════════════════════════════════════════════════════════
%  CHAPTER 7 -- EXPERIMENTAL RESULTS
% ════════════════════════════════════════════════════════════════════════════
\chapter{Experimental Results}

\section{Experimental Setup}

The model was trained on \textbf{75 videos} from SumMe and TVSum.
Key hyperparameters are listed in Table~\ref{tab:hyperparams}.

\begin{table}[H]
\centering
\caption{Model hyperparameters}
\label{tab:hyperparams}
\begin{tabular}{ll}
  \toprule
  \textbf{Parameter} & \textbf{Value} \\
  \midrule
  Input dimension      & 1408 (1024 visual + 384 audio) \\
  LSTM hidden size     & 256 (per direction) \\
  LSTM layers          & 2 \\
  Attention            & Temporal Attention \\
  Optimizer            & Adam \\
  Learning rate        & 0.001 \\
  Frame sampling rate  & 2\,fps \\
  Max sequence length  & 960 frames \\
  Best checkpoint      & Epoch 15 \\
  \bottomrule
\end{tabular}
\end{table}

\section{Quantitative Results}

Table~\ref{tab:quantitative} reports evaluation metrics on the held-out test split.

\begin{table}[H]
\centering
\caption{Quantitative evaluation of the proposed model}
\label{tab:quantitative}
\begin{tabular}{lc}
  \toprule
  \textbf{Metric} & \textbf{Score} \\
  \midrule
  Precision         & 0.4674 \\
  Recall            & 0.4674 \\
  F-score           & 0.4674 \\
  Temporal Overlap  & 0.4222 \\
  \bottomrule
\end{tabular}
\end{table}

\noindent
\textbf{Note:} A direct comparison with published baselines on SumMe/TVSum
\cite{apostolidis2021,zhang2016} is planned as the next step; the above figures
therefore represent preliminary results.

\section{Qualitative Results}

The system was tested on several real-world videos of approximately \textbf{10 minutes},
producing dynamic summaries of \textbf{2--3 minutes} ($\approx 20$--$30\,\%$ of the
original). Key observations:

\begin{itemize}
  \item Selected segments concentrate on regions of high activity, meaningful scene
        transitions, and dialogue containing core semantic content.
  \item Shot-based selection combined with Gaussian smoothing and $\pm 1.5$\,s context
        expansion yields visually coherent, low-jitter summaries.
  \item The retained audio track ensures smooth, uninterrupted sound without abrupt cuts.
  \item The overall viewing experience follows the narrative flow of the original video.
\end{itemize}

% ════════════════════════════════════════════════════════════════════════════
%  CHAPTER 8 -- FUTURE WORK
% ════════════════════════════════════════════════════════════════════════════
\chapter{Future Work}

\section{Extending to Multiple Domains}

\begin{enumerate}
  \item \textbf{Domain-specific video collection:} news broadcasts, educational
        lectures (MOOC/tutorial), and surveillance footage.
  \item \textbf{System evaluation per domain:} apply the current summarization
        pipeline to generate dynamic summaries and record selected keyshot indices
        for downstream use.
\end{enumerate}

\section{Building a Reduced-Dimension Dataset}

The summarized videos serve as a \emph{structurally reduced} version of the original
corpus. This reduced dataset can train downstream models such as:
\begin{itemize}
  \item \textbf{Video classification} (topic labeling of news, lecture type, etc.).
  \item \textbf{Video-to-text} (captioning or question answering on condensed video).
\end{itemize}

\section{Downstream Model Integration}

Candidate architectures for training on summarized video:
\begin{itemize}
  \item \textbf{TimeSformer} -- a Video Vision Transformer with spatiotemporal
        attention, well-suited to shorter, denser sequences produced by summarization.
  \item \textbf{Video Swin Transformer} -- a hierarchical Transformer with 3-D
        shifted windows.
  \item \textbf{I3D / 3D-CNN} -- a simpler baseline for rapid empirical validation.
\end{itemize}

\paragraph{Proposed evaluation protocol.}
\begin{enumerate}
  \item Use the summarization system to reduce the original corpus by $\approx 85\,\%$.
  \item Assign task labels to the summaries.
  \item Train a downstream model on both original and summarized videos.
  \item Compare classification accuracy, training time, and GPU memory usage.
\end{enumerate}

The central hypothesis is that video summarization \textbf{reduces data dimensionality}
while \textbf{preserving} downstream task performance at lower computational cost.

% ════════════════════════════════════════════════════════════════════════════
%  CHAPTER 9 -- CONCLUSION
% ════════════════════════════════════════════════════════════════════════════
\chapter{Conclusion}

This report presented a \textbf{multimodal video summarization system} based on
\textbf{BiLSTM with Temporal Attention} \cite{zhang2016}, combining visual CNN features
with speech-derived semantic embeddings from Whisper \cite{radford2023} and SentenceBERT
\cite{reimers2019}. The system was trained in a supervised fashion on the SumMe
\cite{gygli2014} and TVSum \cite{song2015} benchmarks, achieving a preliminary F-score
of $\approx 0.47$.

Beyond the practical benefit of reducing viewing time, the work demonstrates the role
of video summarization as \textbf{structured dimensionality reduction} -- retaining
information-dense content while discarding redundant frames. This positions the
summarization module as a valuable pre-processing step for downstream video
understanding tasks (classification, retrieval, captioning).

Future directions include a rigorous comparison with published baselines, extension to
news and educational domains, and an empirical study of how summarization-based data
reduction affects the performance and efficiency of modern Transformer-based video
classifiers such as TimeSformer.

% ════════════════════════════════════════════════════════════════════════════
%  BIBLIOGRAPHY
% ════════════════════════════════════════════════════════════════════════════
\bibliographystyle{IEEEtran}

\begin{thebibliography}{11}

\bibitem{haq2021}
H.~B. Haq, M.~Asif, M.~B. Ahmad, R.~Ashraf, and T.~Mahmood,
``Video summarization techniques: A review,''
\textit{Mathematical Problems in Engineering}, vol.~2021, pp.~1--17, 2021.

\bibitem{money2015}
A.~G. Money and H.~Agius,
``A survey on video summarization techniques,''
\textit{International Journal of Computer Applications}, vol.~118, no.~11,
pp.~25--31, 2015.

\bibitem{psallidas2021}
T.~Psallidas, P.~Koromilas, T.~Giannakopoulos, and E.~Spyrou,
``Multimodal summarization of user-generated videos,''
\textit{Applied Sciences}, vol.~11, no.~11, p.~5260, 2021.

\bibitem{hu2023}
S.~Hu, Z.~Liu, J.~Liu, and Z.~Guo,
``Video summarization based on feature fusion and data augmentation,''
\textit{Computers}, vol.~12, no.~9, p.~186, 2023.

\bibitem{kadam2022}
P.~Kadam \textit{et al.},
``Recent challenges and opportunities in video summarization with machine learning
algorithms,''
\textit{IEEE Access}, vol.~10, pp.~122762--122785, 2022.

\bibitem{radford2023}
A.~Radford \textit{et al.},
``Robust speech recognition via large-scale weak supervision,''
in \textit{Proc. 40th Int. Conf. Machine Learning (ICML)}, vol.~202,
pp.~28492--28518, 2023.

\bibitem{apostolidis2021}
E.~Apostolidis, E.~Adamantidou, A.~I. Metsai, V.~Mezaris, and I.~Patras,
``Video summarization using deep neural networks: A survey,''
\textit{Proc. IEEE}, vol.~109, no.~11, pp.~1838--1863, 2021.

\bibitem{gygli2014}
M.~Gygli, H.~Grabner, H.~Riemenschneider, and L.~Van Gool,
``Creating summaries from user videos,''
in \textit{Computer Vision -- ECCV 2014}, Lecture Notes in Computer Science,
vol.~8695, Springer, 2014, pp.~505--520.

\bibitem{zhang2016}
K.~Zhang, W.-L. Chao, F.~Sha, and K.~Grauman,
``Video summarization with long short-term memory,''
in \textit{Computer Vision -- ECCV 2016}, Lecture Notes in Computer Science,
vol.~9911, Springer, 2016, pp.~766--782.

\bibitem{song2015}
Y.~Song, J.~Vallmitjana, A.~Stent, and A.~Jaimes,
``TVSum: Summarizing web videos using titles,''
in \textit{Proc. IEEE Conf. Computer Vision and Pattern Recognition (CVPR)},
2015, pp.~5179--5187.

\bibitem{reimers2019}
N.~Reimers and I.~Gurevych,
``Sentence-BERT: Sentence embeddings using Siamese BERT-networks,''
in \textit{Proc. 2019 Conf. Empirical Methods in Natural Language Processing
(EMNLP-IJCNLP)}, Hong Kong, Nov.~2019, pp.~3982--3992.

\end{thebibliography}

\end{document}